\originalchapter{轻元素核合成与重子密度}
本章对应8.902第19讲末与第20讲。BBN把宇宙膨胀、弱反应和核反应放在同一时间尺度上，能预测氦的大致质量分数，并以氘丰度约束重子密度。解析估计足以看见机制，却不能称作已经执行了完整反应网络的精确丰度计算。下文保留课程使用的近似参数，并指出原笔记把中子平均寿命称为半衰期等会影响结果的区别。

\originalsection{弱反应控制中子与质子的比例}
核子形成后，弱反应$n+\nu_e\rightleftarrows p+e^-$、$n+e^+\rightleftarrows p+\bar\nu_e$以及中子衰变与其逆过程持续改变$n/p$。在充分热平衡、轻子化学势可忽略时，非相对论核子数密度比为
\[
 r_{ \rm eq}\equiv\frac{n_n}{n_p}
 =\left(\frac{m_n}{m_p}\right)^{3/2}
 \exp\left[-\frac{(m_n-m_p)c^2}{\kb T}\right]
 \simeq\exp[-1.293\ \mathrm{MeV}/(\kb T)].
\]
核子自旋自由度相同。温度下降使较重的中子相对减少。弱相互作用的数量级率$\Gamma_w\sim G_F^2T^5$，而辐射时期$H\propto T^2/M_P$，所以$\Gamma_w/H\propto T^3$，最终在MeV量级不再足以维持平衡。

课程取冻结温度约0.8 MeV，瞬时平衡估计$r_f\simeq0.20$。这与中微子逐渐脱耦发生在相近时期，但\textbf{弱反应停止与中微子脱耦不是一个完全相同的瞬间}。真实弱冻结持续一段时间，且需计入中微子温度、电子正电子分布与衰变；后面的氦数字不能由一个固定温度精确确定。

\originalsection{氘瓶颈为何延迟核合成}
第一步$p+n\rightleftarrows\mathrm D+\gamma$形成结合能$B_D\simeq2.22$ MeV的氘。即使$\kb T<B_D$，每个重子对应的光子太多，谱的高能尾部仍可破坏氘。没有足够稳定的氘，快速通向更紧束缚的氦的反应链无法开始，这称氘瓶颈。

在理想非相对论化学平衡中，$\mu_D=\mu_p+\mu_n$。将Maxwell--Boltzmann密度式相除可得
\[
 \frac{n_D}{n_pn_n}=\frac{g_D}{g_pg_n}
 \left(\frac{2\pi\hbar^2}{\kb T}\frac{m_D}{m_pm_n}\right)^{3/2}
 \exp\left(\frac{B_D}{\kb T}\right).
\]
$g_D=3$、$g_p=g_n=2$，忽略核子小质量差而取$m_D\simeq2m_N$，则右侧为$6[\pi\hbar^2/(m_N\kb T)]^{3/2}\exp[B_D/(\kb T)]$。正指数来自氘总质量低于两个自由核子的质量，不能误写成单粒子平衡密度的负Boltzmann指数。

为估算瓶颈解除，仿课程令自由氘数与自由中子数相当，$n_D/n_n\sim1$。这是把约一半尚未进入更重核的中子结合进氘的启发式条件，随后真实核反应网络会立即改变氘丰度。取$n_p\simeq(5/6)\eta_b n_\gamma$、$\eta_b=n_b/n_\gamma$，用黑体光子数公式得到
\begin{equation}\label{eq:astro-deuterium-bottleneck}
 1\sim\frac{10\zeta(3)}{\sqrt\pi}\eta_b
 \left(\frac{\kb T_D}{m_Nc^2}\right)^{3/2}
 \exp\left(\frac{B_D}{\kb T_D}\right).
\end{equation}
系数约6.78，与课程约6.9一致。小$\eta_b\sim10^{-9}$使所需温度远低于$B_D$；反应通常在$\kb T$约0.06--0.1 MeV的范围迅速展开。增加$\eta_b$使同样的结合比例更早、在更高温度达到。

将式\eqref{eq:astro-deuterium-bottleneck}具体取$\eta_b=5\times10^{-10}$、$m_Nc^2=939$ MeV时，可以数值解得约0.066 MeV。电子湮灭后$g_*\simeq3.36$的简化辐射年龄式给约300秒，而较高的0.08 MeV对应约200秒。原课的“$8\times10^8$ K、200秒”是数量级近似，不能与一个特定Saha半结合条件共同声称成精确解；精确时间还依赖自由度演化和非平衡反应。

\originalsection{中子衰变必须同时计入生成的质子}
在弱反应已可忽略、核结合尚少的近似阶段，令$X_n=n_n/(n_n+n_p)$，重子数守恒。中子平均寿命为$\tau_n$，则
\[
 X_n(t)=X_{n,f}\exp[-(t-t_f)/\tau_n],\qquad
 r(t)=\frac{r_f\exp[-\Delta t/\tau_n]}
 {1+r_f[1-\exp(-\Delta t/\tau_n)]}.
\]
分母增加来自衰变产生的质子。只有$r_f$很小时才近似$r\simeq r_f\exp(-\Delta t/\tau_n)$。课程所用890秒应理解为\emph{平均寿命}；其对应半衰期为$\tau_n\ln2\simeq617$秒。若真的把890秒当成半衰期，应使用$2^{-\Delta t/890\mathrm s}$，会给不同答案。

例如$r_f=0.20$、$\Delta t=200$秒、$\tau_n=890$秒给$r\simeq0.154$，并非由此精确推出0.25的氦分数。有效冻结温度、瓶颈结束时刻和残余弱反应的改进均重要，原讲义把后面的步骤称“精确计算”超过了这些近似所能支持的精度。本卷采用它说明数量级和参数方向，不伪造一个没有运行的精确网络结果。

\originalsection{氦质量分数的简洁估计}
瓶颈解除后，氘、氚与${}^3\mathrm{He}$反应迅速将大部分剩余中子锁入强束缚的${}^4\mathrm{He}$。每个氦核需两个中子，若忽略少量其他核和结合能对质量的修正，有$n_{\rm He}=n_n/2$、$n_{\rm H}=n_p-n_n$，因此
\begin{equation}\label{eq:astro-helium-fraction}
 Y_p=\frac{4n_{\rm He}}{n_{ \rm H}+4n_{\rm He}}
 =\frac{2n_n}{n_p+n_n}=2X_n=\frac{2r}{1+r}.
\end{equation}
若$r=1/7$，便得$Y_p=1/4$。若所有冻结时$r_f=1/5$的中子瞬间结合，则给$1/3$；延迟期间的中子衰变使结果下降。上面200秒例的$Y_p\simeq0.266$直观表明，只改变最后一位四舍五入不能把近似模型变成精确预测。

大部分核合成停止后，仍留下少量氘、${}^3\mathrm{He}$和${}^7\mathrm{Li}$或随后衰变成锂的${}^7\mathrm{Be}$。缺乏稳定的质量数5、8核，加上密度迅速下降与Coulomb势垒，使普通BBN难以生产大量更重元素；恒星内部的高密度、长时间和三氦过程才有不同机会。不能由本章一个氦式推算全部轻元素丰度。

\originalsection{从丰度到重子参数，需要哪些额外信息}
较高$\eta_b$提高核反应相对于膨胀的速率，让氘更有效地被烧成氦，故剩余D/H通常下降；氦分数对$\eta_b$依赖较弱。额外辐射自由度提高$H$，弱反应更早冻结且核合成前时间缩短，通常增加中子存留和氦分数。中子寿命、轻子化学势与反应截面也进入定量推断，不能将$Y_p$只看成$\Omega_b$的单参量函数。

对既定今日CMB温度，由$n_{b,0}=\Omega_b\rho_{ \rm crit,0}/m_p$与黑体$n_{\gamma,0}$得
\[
 \eta_b\simeq2.74\times10^{-8}\,\Omega_bh^2
 \left(\frac{2.7255\ \mathrm K}{T_0}\right)^3.
\]
因此氘等轻元素约束的是物理重子密度$\omega_b=\Omega_bh^2$；若要单独给$\Omega_b$，还需$h$。将BBN与CMB的重子加载比较，是不同物理过程的交叉核验，不是用同一份观测重新作两次拟合。

观测原初氦需从低金属丰度区域外推并控制恒星污染，氘常用远处吸收系统中的同位素谱线，需排除氢线混叠和化学演化。标准核网络与贫金属恒星锂丰度的差异是需要额外物理或天体系统误差解释的难题，本卷不凭一段解析估计声称已解决。反应网络、误差传播及真实数据拟合未实施，范围在覆盖清单单列。
